%======================================================================
\chapter{Proposal}
%======================================================================

% The proposal is a document that outlines the topic of your project. Effectively, the proposal should describe why the
% scope of the project is appropriate for this course. The course instructor will provide feedback on the proposal
% following submission. Students are encouraged to be as detailed as possible in the proposal. The more details
% provided in the proposal, the more detailed feedback you will receive.

% The proposal should contain the following elements:
% 1. Background and objective(s) for the project.
% 2. Proposed method(s) from artificial intelligence.
% 3. Dataset or environment to be used.
% 4. Proposed validation/analysis strategy.
% 5. Description of the novelty of the project.
% 6. Schedule of milestones to be completed each week (beginning the week after the proposal is submitted).

% Page limit for proposals: 2 pages maximum
% Please use a standard page format (i.e. page size 8.5” x 11”, minimum ½” margins, minimum font size 10, no
% condensed fonts or spacing). Please submit the proposal as a PDF file.
% Proposals are to be submitted electronically through Brightspace. It is your responsibility to ensure that your
% proposal is submitted properly. Students may submit the proposal earlier to receive earlier feedback on the proposal.
% The proposal is worth 12% of the project grade.

\section{Background and Objective}

For decades, we've been trying to completely erase the language barrier in communication. As advancemants have been made in machine learning and NLP, we've recently seen the rise of \textit{speech-to-speech} translation, where AI-generated voices could "dub" a speaker in real time. This is a huge step forward, but while these milestones have made the words understandable, they have failed to solve the \textit{"Immersion Gap"}.

\subsection{The Visual Frontier}

In a natural human-to-human conversation, \textit{nearly 70\%} of communication is non-verbal \cite{mehrabian1971silent}; the visual medium holds the key to truly erasing the digital language barrier. Today, even the most advanced video calls suffer from a jarring "visual disconnect." When a speaker's mouth moves to the rhythm of Spanish while the listener hears a translated English stream, the brain flags this as uncanny and distracts listeners from your message. We have solved the problem of what is being said, but we have yet to solve how it is seen.

\subsection{The Problem}

Existing solutions for real-time lip-synchronization (like Wav2Lip, FlawlessAI, or diffusion-based models) currently face issues including but not limited to the following:

\begin{enumerate}
    \item \textbf{Production vs. Live Utility}: Many existing systems are strictly tethered to offline post-production rather than the spontaneity of a live video call
    \item \textbf{High Latency}: Produces "sync-lag" that disrupts the flow of live dialogue
    \item \textbf{Proprietary \& Rigid}: Often tied to expensive software licenses or closed ecosystems
\end{enumerate}

\subsection{Our Objective}

The objective of this project is to bridge the final gap in cross-lingual communication by shipping a \textbf{Real-Time, Lightweight Lip-Sync Engine} specifically optimized for online video conferencing. It should be able to attain the following milestones:

\begin{enumerate}
    \item \textbf{Core MVP Development}: Build an end-to-end pipeline that ingests live audio and outputs a temporally consistent, lip-synced video stream for real-time translation
    \item \textbf{Real-Time Performance}: Achieve accurate lip-syncing with a latency of under 500ms on common consumer-grade hardware
    \item \textbf{Cross-Platform Compatibility}: Ensure compatibility with major video conferencing platforms (Zoom, MSTeams, Google Meet)
\end{enumerate}

\section{Methods}

To achieve low-latency, photorealistic lip-syncing for live video conferencing, this project utilizes an Audio-to-Blendshape translation pipeline. The system maps incoming translated audio to physical parameters on a pre-registered 3D facial mesh, completely bypassing the visual distortion and high computational cost of 2D pixel generation.

\subsection{3D Facial Model Acquisition}
Before any inference occurs, a static 3D digital model of the user's head must be acquired. We will utlilze a camera with TrueDepth or an equivalent feature and a 3D scanning application to capture a high-fidelity topological mesh of the user's head. The mesh will then be rigged with a standard 52 ARKit facial blendshapes. This pre-registered, rigged 3D model will what our neural network animates.
\subsection{Audio Translation Pipeline}
To generate translated audio, we will implement a three-stage pipeline to convert the speaker's voice into a translated audio stream in the target language using lightweight models:
\begin{enumerate}
    \item \textbf{Stage 1:} Speech-to-Text using an Automatic Speech Recognition (ASR) model in the original language
    \item \textbf{Stage 2:} Machine Translation (MT) to translate the text to the target language
    \item \textbf{Stage 3:} Text-to-Speech vocoder to generate the translated audio stream in the target language
\end{enumerate}

\subsection{Visual Lip-Sync Pipeline}
The core of this project is revolves around the visual lip-sync pipeline, which takes the translated audio stream and generates a lip-synced video output. The pipeline consists of the following stages:
\begin{enumerate} 
    \item \textbf{Data Representation:} The synthesized audio is converted into Mel-Frequency Cepstrum Coefficients (MFCC) to extract phonetic features over a sliding time window.
    \item \textbf{Model Architecture:} We will implement a hybrid 1D-CNN and LSTM architecture.
    \begin{itemize}
        \item \textbf{Feature Extraction:} The 1D-CNN layer will process the MFCCs and extract local phonetic features.
        \item \textbf{Temporal Context:} The LSTM layer processes the sequence to maintain temporal context.
        \item \textbf{Output Layer:} A dense layer that maps the LSTM output to the 52 blendshape coefficients, which will be used to animate the 3D facial model.
    \end{itemize}
\end{enumerate}

\subsection{3D Rendering and Overlay}
During live inference, the network's predicted 52 values are passed to a lightweight 3D rendering script. This script will animate the pre-registered 3D mesh acquired in the first step. Then using standard CV techniques like OpenCV, the script will extract the animated 3D mouth/jaw region and seamlessly overlay it onto the user's face in the live 2D webcam feed.

\section{Dataset}

In order to train the audio-to-blendshape neural networks, we will utilize the BEAT dataset (Body-Expression of Audio-Driven Talking) \cite{liu2022beat}, specifically using the 4-modality captured data.

\subsection{Dataset Identification and Composition}
The BEAT dataset is a large-scale, high quality, multi-modal dataset contianing over 76 hours of audio-visual data. Crucially for our translation pipeline, the dataset is multi-lingual, comprising of conversational speech in English (60h), Chinese (12h), Spanish (2h), and Japanese (2h). The dataset natively provides raw audio tracks paired precisely to 52 ARKit facial blendshape weights captured at 60 fps. 
\subsection{Task Environment and Procurement Protocol}
The dataset will be procured via the official academic repository. To prepare the environment for our neural network, we will implement the following data pipeline:
\begin{enumerate}
    \item \textbf{Audio Preprocessing:} The multi-lingual audio files will be uniformly resampled and converted into sequential MFCCs using the librosa library. These MFCCs will serve as the continuous input features $X$. 
    \item \textbf{Target Extraction:} The ground-truth ARKit blendshape weights, ranging from 0.0 to 1.0, will be extracted from the dataset and convereted into PyTorch tensors. These tensors will serve as our continuous target variables $Y$. 
    \item \textbf{Data Splitting:} The full dataset will be randomized and split into 80/10/10 training, validation, test splits.
\end{enumerate} 

\section{Validation Strategy}

To evalute the performance of our product and ensure the lip-syncing is suitable for a live vide-conferencing environment, we will perform a combination of quantitative and qualitative evaluations:
\subsection{Quantitative Validation:} The primary objective of the LSTM is to accurately predict the physical geometry of the mouth. We will evaluate the performance using MSE on the predicted blendshape vectors and the held-out test sets ground-truth blendshape vectors. We may also calculate the velocity of predicted blendshapes over time compared to ground truth to quantify and penalize unrealizstic, jittery movements. We will also measure the end-to-end latency of the system to ensure it meets our real-time performance target.
\subsection{Qualitative Validation:} We will perform eye tests to evaluate certain phonetic mouth shapes being animated properly, and conduct user studies to evaluate the overall visual quality and immersion of the lip-sync output. 

\section{Novelty of Project}

The current market for lip-sync technology is polarized. 

On one end are enterprise-grade solutions like FlawlessAI, which deliver cinema-quality results but are strictly tethered to asynchronous post-production and massive GPU clusters. They are tools for studios with proprietary licensing, not people. On the other end are lightweight, open-source models that run on consumer hardware but produce "deepfakes" so uncanny and distorted they become a distraction rather than a communication aid.

The novelty of our project ultimately lies in balancing real-time performance with a lightweight architecture that brings high-fidelity lip-syncing to standard consumer hardware. Additionally, our focus on cross-platform compatibility ensures that our solution can be easily integrated into popular video conferencing platforms, further enhancing its usability and impact.

\section{Timeline and Milestones}

\begin{itemize}
    \item \textbf{Week 1 (February 23 - March 1):} Finalize real-time audio pipeline, acquire 3D data, and begin model architecture experiments.
    \item \textbf{Week 2 (March 2 - March 8):} Implement and train the audio-to-blendshape model and evaluate performance for latency bottlenecks.
    \item \textbf{Week 3 (March 9 - March 15):} Complete MVP by integrating the trained model with 3D rendering and conducting integration tests.
    \item \textbf{Week 4 (March 16 - March 22):} Optimize pipeline for fidelity, implement cross-platform features, and perform end-to-end testing.
    \item \textbf{Week 5 (March 23 - March 29):} Document all project results and prepare the final report and presentation materials.
\end{itemize}